\section{Cross-reference graph}
\label{sec:graph}

With passages labelled down to their subdivisions, the references between them
can be resolved. The graph is built once per upload, stored beside the vectors,
and read at answer time. This section follows it through its five stages:
extracting references, turning them into edges, storing them, expanding the
neighbourhood of an answer, and presenting it to the reader.

\subsection{Extracting references}

A reference is a unit word followed by an identifier and, optionally, a
paragraph in parentheses, as in ``Article 5(2)''. The extractor uses the same
unit vocabulary as the structure parser, so anything that can be cited can
also be a heading, and the reverse. A match that stands alone on its line is a
heading, not a reference, and is skipped.

Each reference is then classified as internal or external. It points at
another instrument when it is followed by ``of'' and a capitalised name that
is not itself a unit (``Article 11 of Regulation (EU) 2019/2144''), or by an
acronym of three or more capitals (``Article 114 TFEU''). ``Article 5 of this
Regulation'' stays internal, because ``this'' is not a name. References are
deduplicated per passage by unit, identifier and instrument.
\Cref{lst:references} shows the result on a sample passage.

\begin{lstlisting}[float, numbers=none, caption={References extracted from a sample
passage.}, label={lst:references}]
1. As referred to in Article 5(2) and Annex III, the requirements of
Chapter III, Section 2 apply. Article 11 of Regulation (EU) 2019/2144
applies, as does Article 114 TFEU. See also Article 5 of this Regulation.

article 5, paragraph 2    internal
annex III                 internal
chapter III               internal
section 2                 internal
article 11                external: Regulation (EU) 2019/2144
article 114               external: TFEU
\end{lstlisting}

\subsection{Building edges}

A node is a provision: the most specific provision-level label a passage
carries---recital, article, clause, rule, section, annex, schedule, appendix,
exhibit, chapter, part or title---written as \texttt{kind:value}, for example
\texttt{article:66}. The builder reads the document's passages back from the
vector store in reading order and, for each reference in each passage, adds an
edge from the passage's provision to the provision cited. An edge is
\emph{internal} when some passage of the document carries the cited label,
\emph{external} when it names another instrument, and \emph{unresolved} when
the document has no passage with that label. A provision's references to
itself are skipped. Every edge also records its \emph{position}, the index of
the citing passage, which orders a provision's references as they appear in
the text, and its \emph{locator}, the paragraph, subparagraph, point and
sub-point of the citing passage (``paragraph 1 $\cdot$ point b'').

Reading the passages back from the store, rather than from the parser's
output, guarantees that edges use exactly the labels retrieval uses: a node
identifier converts directly into the metadata filter that fetches its text.

\paragraph{Trade-offs.} Nodes are provisions, not paragraphs. References to
``Article 5(2)'' and ``Article 5(3)'' therefore reach the same node, and the
precision on the cited side is traded for a graph that stays small and whose
every node can be fetched with one filter. Precision on the citing side is
kept in the locator. For the AI Act the graph holds 575 internal edges, 79
edges to other instruments and 29 unresolved ones, the latter mostly
references to units the parser does not label as provisions, such as the
sections inside an annex.

\subsection{Storage}

Edges are stored in SQLite (\cref{lst:schema}). A document's edges are
replaced whenever it is uploaded and removed when it is deleted. Two queries
serve every use: the edges leaving a set of nodes, ordered by position, and
the edges arriving at a set of nodes, which the index on the target makes
cheap.

\begin{lstlisting}[float, language=SQL, caption={The edge table
(\texttt{graph/store.py}).}, label={lst:schema}]
CREATE TABLE IF NOT EXISTS edges (
    position    INTEGER NOT NULL,
    locator     TEXT NOT NULL,
    document    TEXT NOT NULL,
    source      TEXT NOT NULL,
    target      TEXT NOT NULL,
    edge_type   TEXT NOT NULL,
    PRIMARY KEY (document, source, target, position)
);
CREATE INDEX IF NOT EXISTS edges_by_target ON edges (document, target);
\end{lstlisting}

\paragraph{Trade-offs.} A graph database would add a server to install and
run, for queries that are all one hop long. SQLite ships with Python, keeps
the edges in a single file beside the vector store, and answers both queries
from an index. Traversals longer than one hop are done in Python with a
bounded depth, which is all the interface needs.

\subsection{Neighbourhood of an answer}

Every answer carries the cross-reference neighbourhood of its passages. The
provisions of the retrieved passages are the seeds, at depth 0. The
neighbourhood is then expanded breadth-first along internal edges to depth 2,
with at most 50 nodes. Each node carries its label, its kind (retrieved,
internal, external or unresolved), its depth and, for provisions in the
document, a preview of its first 240 characters. Each edge carries its type,
its locator and an index that numbers a provision's references in reading
order. Because the first level arrives with the answer, the references of a
retrieved passage can be listed without another request.

Two endpoints serve deeper exploration on demand: one returns what a
provision cites and what cites it, each in reading order, and the other
returns a provision's full text (\cref{tab:api}).

\subsection{Presenting cross-references}

\paragraph{A graph view, and why it was dropped.} The first presentation was
an interactive graph drawn beside the answer, first with a force-directed
layout and then with a radial one that placed the retrieved provisions at the
centre, arranged each further level on a wider ring, swapped nodes to shorten
and uncross the arrows, and routed curved arrows around the boxes. It took
about 1\,500 lines of TypeScript. With the typical neighbourhood of 20 to 50
nodes it showed \emph{that} provisions referred to each other, but not
\emph{what} they said: reading a chain meant opening node after node while
holding the path in memory, and the layout's motion competed with the text.
The graph view was therefore dropped in favour of a much simpler interface,
and the graph itself remained in the backend as the data behind it.

\paragraph{The reference tree.} Under each citation card, the provisions the
passage cites are listed in the order the text cites them, each with the
locator of the citing passage (\cref{lst:tree}). Opening a row fetches the
cited provision's text and its own references once and shows them indented
beneath it, so a chain of citations stays on screen and reads from top to
bottom. A reference to another instrument is marked \emph{other document}, an
unresolved one \emph{not found}, and a provision that already appears higher
on the same path \emph{already above}; none of these can be opened, which also
makes every chain finite even when provisions cite each other. A long list
shows twelve rows and offers the rest on request. Nothing animates, and every
row exposes its open or closed state to assistive technology.

\begin{lstlisting}[float, numbers=none, caption={The reference tree under Article
101(1), with two levels opened (\texttt{v} open, \texttt{>} closed).},
label={lst:tree}]
Cites
v 1. Article 91 . paragraph 1 . point b
     Article 91 Power to request documentation and information
     1. The Commission may request the provider of the general-purpose
        AI model concerned to provide the documentation ...
     v 1. Article 68 . paragraph 3
          Article 68 Scientific panel of independent experts
          1. The Commission shall, by means of an implementing act, ...
> 2. Article 93 . paragraph 1 . point c
\end{lstlisting}

\paragraph{Trade-offs.} The tree gives up the overview of the whole
neighbourhood that a graph offers at a glance, and with it the ability to see
clusters and hubs. In exchange, the text of every cited provision is readable
in place and in order, which is what following a cross-reference is for. The
component is about 140 lines, a tenth of the graph view, and costs two small
requests per opened row instead of a layout computation on every change.
